\section{Limitations and Broader Impact}

\method{} assumes that augmentations preserve class identity and that the source
and target label spaces coincide. It may therefore fail under label shift,
open-set targets, or fine-grained tasks for which strong augmentation removes
class-defining evidence. The online thresholds describe model confidence, not
ground-truth correctness; systematic, augmentation-stable errors can still pass
both gates. Future work should combine the gate with explicit open-set detection
and study guarantees under target-label shift.

More reliable adaptation can reduce annotation requirements, but it can also
make deployed models appear trustworthy despite unobserved failure modes.
Practitioners should report per-group performance where appropriate, monitor
coverage drift, and retain human review in consequential applications. We do
not recommend interpreting confidence gating as a substitute for external
validation on representative target data.

% -----------------------------------------------------------------------------
